\documentclass[11pt]{article}
\usepackage[T1]{fontenc}
\usepackage{lmodern}
\usepackage[letterpaper,margin=0.8in]{geometry}
\usepackage{booktabs,multirow,pdflscape,rotating}
\setlength{\parindent}{0pt}
\setlength{\parskip}{8pt}
\renewcommand{\arraystretch}{1.4}
\setlength{\tabcolsep}{8pt}
\begin{document}
\begin{landscape}
\section*{1. Landscape page}
The page displays horizontally in a PDF viewer. The table remains upright
relative to the displayed page. This gives a wide table more horizontal space.
\begin{center}
\small <<revenue_table>>
\end{center}
Revenue in USD thousands. Illustrative data.
\end{landscape}

\begin{sidewaystable}[p]
\centering
\begin{minipage}{\textheight}
\section*{2. Sideways table on a portrait page}
The PDF page stays portrait. This entire table block is rotated 90 degrees,
including its title and explanatory text. Rotate the viewer to read it upright.
\begin{center}
\small <<revenue_table>>
\end{center}
Revenue in USD thousands. Illustrative data.
\end{minipage}
\end{sidewaystable}
\end{document}
